\section{Experiments}
\label{sec:experiments}

In this section, numerical experiments are presented to show the performance of the proposed DM algorithm in preserving the geometric properties of the underlying manifold, particularly focusing on reach and Riemannian metric. Manifolds are generated by the package in \cite{JMLR:v23:19-644}. For details of generating data, please refer to Appendix \ref{more_exp}. All experiments were conducted with $10000$ sample points on a MacBook Air with an Apple M4 processor and standard macOS system configuration, the running time is around $30$ minutes for each experiment. 

\subsection{Common Manifolds}

The first experiment is testing the reach preservation and Riemannian metric (Rmetric) preservation on synthetic datasets sampled from known manifolds. 
The process is the following: 
\begin{enumerate}
    \item Adopting the method in \cite{rmetric,joncas2017improved}, choosing the best $h$ in \cref{Lap matrix} and neighborhood radius $r$ using grid search based on empirical validation for estimating Rmetric. In this step, Rmetric (baseline) and reach are estimated. If the manifold has a boundary, then we only consider estimation at the inner points~(Appendix \ref{more_exp}) to avoid the bias caused by boundary points.
    \item Selecting best $t$ for the DM algorithm according to \cite{rmetric,joncas2017improved}, and embed the data into $\R^3$, then estimate the reach and pushforward Riemannian metric for the embedded manifold.
\end{enumerate}

$\text{Rmetirc--error}_{\text{baseline}}$ ($\text{R--error}_{\text{base}}$) is computed directly on the original manifold and measures the intrinsic estimation error of the metric estimation procedure. Since the ground-truth pushforward metric is equal to $1$, the baseline error is defined as the deviation of the estimated metric from $1$.
$\text{R--error}_{\text{emb}}$ is computed on the embedded manifold and measures how well the DM embedding preserves the Riemannian metric. It is defined as the relative error with respect to the baseline estimate.
The reach is estimated using definitions from \cite{aamari2019estimating}. The table below summarizes the results.


\begin{table}[htbp]
    \caption{Results of the experiments on synthetic datasets.}
    \centering
    \begin{tabular}{cccccccc}
    \toprule
    Manifolds & $\text{R--error}_{\text{base}}$ & $\text{reach}$ & $\text{R--error}_{\text{emb}}$ & $\text{reach}_{\text{emb}}$ & $h$ & $r$ & $t$ \\
    \midrule
    Sphere & 0.0125 & 0.9999 & 0.0203 & 0.9441 & 0.5198 & 1.578 & 0.3360\\
    Swiss Roll & 0.1085 & 0.4468 & 0.0933 & 0.3459 & 0.9161 & 2.689 & 1.880 \\
    Torus & 0.1564 & 0.3162 & 0.1562 & 0.1731 & 0.9190 & 2.812 & 1.788\\
    \bottomrule
    \end{tabular}
\end{table}


The first manifold is $S^2$ since $S^2$ is well-studied and has analytical eigenvalues and eigenfunctions of the Laplace--Beltrami operator, we can compare the estimated results with true values easily. The second and third manifolds are Swiss Roll and Torus (Standardized), respectively, which are more complex manifolds with non-trivial geometry.

The results show that the DM algorithm preserves the reach and Riemannian metric well for the sphere, with only a slight increase in R--error and a small decrease in reach. For the Swiss Roll and Torus, original reach is accurately estimated, where the true values of reach are approximately $\pi/6.97 = 0.4507$ and $0.668^2/1.409 = 0.3167$ (constants are related to standardization scaling factors), respectively. However, the Rmetric cannot be captured as well as the sphere, since these manifolds have more complex geometry and the algorithm cannot handle them effectively, where larger sample sizes are needed to decrease the error. And an increase in R--error on the original manifold is associated with a corresponding increase in R--error in the embedded manifold, along with a larger deviation in reach.

\subsection{Random Manifolds}
The DM algorithm is also tested on 2 randomly generated complex manifolds (for more details see appendix \ref{more_exp}), and the results are summarized in the table below.


\begin{table}[htbp]
    \caption{Results of the experiments on randomly generated manifolds, where the baseline is the R--error computed on the original manifold, and other columns are the relative R--error computed on the embedded manifold with different embedding dimensions.}
    \centering
    \begin{tabular}{lccccccc}
    \toprule
    \diagbox{Manifold}{Dim} & baseline & $10$ & $20$ & $30$ & $40$ & $h$ & $r$\\
    \midrule
    Manifold 1 & 0.2312 & 0.5528 & 0.4272 & 0.3530 & 0.3276 & 0.1718 & 0.3938\\
    Manifold 2 & 0.3286 & 0.5945 & 0.3909 & 0.2797 & 0.2440 & 0.3006 & 0.6789\\
    \bottomrule
    \end{tabular}
\end{table}


The results suggest that as the embedding dimension increases, the R--error decreases, indicating better preservation of the Riemannian metric. And the decreasing rate is slower when the embedding dimension is larger, which is expected since the eigenvalues increase exponentially, and the contribution of higher order eigenvectors to the embedding is weakened by the eigenvalues. 

However, even when the embedding dimension is $40$, the R--error remains relatively large. One possible reason is that the algorithm for estimating Riemannian metric cannot perform well on complex manifolds (as presented by the baseline error), so that even almost isometric embeddings can have large error when estimating from point clouds.